% Department presentation for junior students. Overleaf: Main document = presentation.tex.
% Data: data/site/ (scripts/sync_site.py), data/research/ (scripts/research_stats.py).

\documentclass[aspectratio=169, 9pt]{beamer}
\usetheme{intsystems}

\title[Кафедра интеллектуальных систем МФТИ \quad intsystems.github.io]{Кафедра интеллектуальных систем}
\date{\today}

\begin{document}

\coverframe{%
  \covertitle[1.2cm]{Кафедра\\интеллектуальных\\систем}{Физтех-школа прикладной математики и~информатики\\МФТИ}%
}

\begin{frame}{Кафедра в цифрах}
  \vspace{0.5cm}
  \begin{columns}[T]
    \begin{column}{0.3\textwidth}\statvalue{2003}{год основания кафедры}\end{column}
    \begin{column}{0.3\textwidth}\statvalue{\research{alumni}}{выпускников с~\research{alumni-from} года}\end{column}
    \begin{column}{0.3\textwidth}\statvalue{\thesestotal}{дипломов защищено за 2022--2026}\end{column}
  \end{columns}
  \vspace{0.9cm}
  \begin{columns}[T]
    \begin{column}{0.3\textwidth}\statvalue{\githubtotal}{открытых проектов на GitHub}\end{column}
    \begin{column}{0.3\textwidth}\statvalue{<35}{средний возраст преподавателей курсов}\end{column}
    \begin{column}{0.3\textwidth}\statvalue{>50\,\%}{выпускников защитили кандидатские}\end{column}
  \end{columns}
  \vfill
\end{frame}

\begin{frame}{О кафедре}
  \begin{columns}[T]
    \begin{column}{0.5\textwidth}
      \subhead{Направления}
      \begin{itemize}
        \item Бакалавриат --- 03.03.01 «Прикладные математика и физика»
        \item Магистратура --- 09.04.01 «Информатика и вычислительная техника»
        \item Аспирантура --- 1.2.1 «Искусственный интеллект и машинное обучение»
      \end{itemize}
      \subhead{Базовая организация}
      Вычислительный центр им.~А.\,А.~Дородницына ФИЦ «Информатика и~управление» РАН
      \vspace{0.25cm}

      \includegraphics[height=0.5cm]{fpmi.png}\hspace{0.8cm}%
      \includegraphics[height=0.9cm]{vc-ras.png}
    \end{column}
    \begin{column}{0.46\textwidth}
      \card{\subhead{Теория}
        \footnotesize Машинное обучение, глубокое обучение, генеративные нейросети, байесовские модели и~мультимодели, функциональный анализ данных}
      \vspace{0.2cm}

      \card{\subhead{Практика}
        \footnotesize Геометрические глубокие нейросети, анализ текстов и тематическое моделирование, анализ изображений, видео и~аудио, биомедицинские сигналы, интерфейсы человек--компьютер}
    \end{column}
  \end{columns}
\end{frame}

\begin{frame}{История}
  \begin{columns}[T]
    \begin{column}{0.48\textwidth}
      \card[\linewidth][2.05cm]{\subhead{Интеллектуальный анализ данных}
        Создана в 2003 году академиком РАН К.\,В.~Рудаковым на базе одной из старейших кафедр ФУПМ МФТИ «Управление и~вычислительные системы».}
    \end{column}
    \begin{column}{0.48\textwidth}
      \card[\linewidth][2.05cm]{\subhead{Проектирование и организация систем}
        Создана в 1970-е годы академиком Г.\,С.~Поспеловым, основоположником отечественной школы искусственного интеллекта, как базовая кафедра при ВЦ РАН.}
    \end{column}
  \end{columns}
  \vfill
  \centering
  \foreach \file/\name in {zhuravlyov_yv/Ю.\,И.~Журавлёв, rudakov_kv/К.\,В.~Рудаков, vorontsov_kv/К.\,В.~Воронцов, pospelov_gs/Г.\,С.~Поспелов} {%
    \begin{tabular}{@{}c@{}}
      \cropphoto{\file.jpg}{1.75cm}{2.3cm}\\[0.1cm]
      \scriptsize\name
    \end{tabular}\hspace{0.6cm}%
  }
  \vfill
\end{frame}

\begin{frame}{Студенты кафедры}
  \begin{columns}[T]
    \begin{column}{0.66\textwidth}
      \begin{itemize}
        \item Ежегодно кафедра принимает и выпускает \semibold{10--15 студентов} в~бакалавриат и магистратуру
        \item Одновременно на кафедре учатся \semibold{до 50 студентов}
        \item Часть студентов параллельно учится на международных программах: \semibold{Университет Гренобль-Альпы}, \semibold{Политехническая школа} (Париж), \semibold{KAUST}
        \item После защиты выпускники идут в аспирантуру МФТИ и~ведущих университетов, в~исследовательские лаборатории и~компании
      \end{itemize}
    \end{column}
    \begin{column}{0.3\textwidth}
      \centering
      \includegraphics[width=0.7\linewidth]{uga.jpg}\\[0.25cm]
      \includegraphics[width=0.38\linewidth]{polytechnique.png}\\[0.25cm]
      \includegraphics[width=0.7\linewidth]{kaust.png}
    \end{column}
  \end{columns}
\end{frame}

\begin{frame}{Исследования в центре учёбы}
  \begin{columns}[T]
    \begin{column}{0.5\textwidth}
      \researchsteps{%
        3 курс/собеседование и выбор научного руководителя,
        Моя первая научная статья/первая статья за семестр: paper{,} code{,} slides,
        Каждый семестр/отчёт о НИР: статья{,} код эксперимента{,} доклад,
        Конференции/NeurIPS{,} ICML{,} ICLR{,} AISTATS{,} конференция МФТИ,
        Защита/диплом с работой{,} поданной к публикации}
    \end{column}
    \begin{column}{0.46\textwidth}
      \subhead{Защиты дипломов по годам}
      \thesischart
    \end{column}
  \end{columns}
\end{frame}

\begin{frame}{Бакалавриат: курсы кафедры}
  \coursecards[2]{bachelor}
  \vfill
  \muted{Итог бакалавриата --- статья и код проекта; их учитывают при поступлении на «магистратуру плюс аспирантуру» в МФТИ и~других университетах.}
\end{frame}

\begin{frame}{Магистратура: курсы кафедры}
  \coursecards[3]{master}
  \vfill
  \muted{Выпускники магистратуры --- авторы нескольких статей: этого достаточно для поступления в~аспирантуру МФТИ и~ведущих университетов.}
\end{frame}

\begin{frame}{Преподаватели}
  \vspace{0.2cm}
  \teamgrid{\instructors}
\end{frame}

\begin{frame}{Открытый код}
  \begin{columns}[T]
    \begin{column}{0.72\textwidth}
      \subhead{Новые открытые проекты по годам: всего \githubtotal}
      \githubchart[4.6cm]
    \end{column}
    \begin{column}{0.24\textwidth}
      \centering
      \vspace{0.6cm}
      \qrlink[2cm]{https://github.com/intsystems}{github.com/intsystems}\\[0.4cm]
      \qrlink[2cm]{https://www.youtube.com/c/MachineLearningPhystech}{YouTube: лекции и доклады}
    \end{column}
  \end{columns}
\end{frame}

\begin{frame}{Выпускники в науке}
  \ifresearch{alumni-matched}{%
    \begin{columns}[T]
      \begin{column}{0.5\textwidth}
        \cohortchart[2.5cm]{cites}{медиана цитирований}\par\vspace{0.1cm}
        \cohortchart[2.5cm]{h}{медиана индекса Хирша}
      \end{column}
      \begin{column}{0.46\textwidth}
        \subhead{Самые цитируемые выпускники}
        \scriptsize\topalumni
      \end{column}
    \end{columns}
  }{\card{Данные OpenAlex не загружены: запустите scripts/research\_stats.py.}}
\end{frame}

\begin{frame}{Выпускники последних пяти лет}
  \ifresearch{alumni-matched}{%
    \subhead{Самые цитируемые выпускники \the\numexpr\year-5\relax--\the\year{} годов}
    \vspace{0.2cm}\small\topalumni[recent][0.7em]
  }{\card{Данные OpenAlex не загружены: запустите scripts/research\_stats.py.}}
\end{frame}

\begin{frame}{Где работают выпускники}
  \begin{columns}[T]
    \begin{column}{0.48\textwidth}
      \ifresearch{alumni-placed}{%
        \subhead{Последняя аффилиация в публикациях}
        \sectorbars
      }{\card{Данные OpenAlex не загружены: запустите scripts/research\_stats.py.}}
      \subhead{Стажировки и совместные проекты}
      \foreach \logo in {antiplagiat.jpg, forecsys.png, goods-forecast.png, yandex.png, sber.jpg, avito.png, huawei.png, ozon.png, mipt-mil.png} {%
        \makebox[0.32\linewidth]{\includegraphics[height=0.6cm, width=0.26\linewidth, keepaspectratio]{\logo}}\hspace{0pt}%
      }
    \end{column}
    \begin{column}{0.48\textwidth}
      \ifresearch{abroad}{%
        \subhead{Зарубежные организации выпускников}
        \card{\footnotesize\research{abroad}}
      }{}
    \end{column}
  \end{columns}
\end{frame}

\begin{frame}{Научная стипендия им.~К.\,В.~Рудакова}
  \begin{columns}[T]
    \begin{column}{0.72\textwidth}
      \begin{itemize}
        \item Для бакалавров и магистров кафедры, назначается на один семестр
        \item \semibold{20\,000 рублей} в месяц
        \item Заявки: до 2 сентября и до 15 февраля, intsystems.github.io/materials/scholarship
      \end{itemize}
      \subhead{Стипендиаты осени 2026 года}
      \scriptsize
      \begin{enumerate}
        \item \semibold{Алексей Кравацкий} --- ImprovEvolve: Ask AlphaEvolve to Improve the Input Solution and Then Improvise
        \item \semibold{Ксения Ибрагимова}
        \item \semibold{Дмитрий Василенко} --- Jacobian Analysis of a Recurrent Transformer Block
        \item \semibold{Владислав Мешков} --- Model Compression via Internal Regularization: A Comparative Study of Sparsity-Inducing Algorithms for Layer Pruning
        \item \semibold{Фёдор Соболевский} --- Automatic Hierarchical Summarization of Scientific Texts
      \end{enumerate}
    \end{column}
    \begin{column}{0.24\textwidth}
      \raggedleft\roundphoto{rudakov_kv.jpg}{2.4cm}
    \end{column}
  \end{columns}
\end{frame}

\begin{frame}{Конференции}
  Сотрудники и студенты кафедры публикуются в высокорейтинговых журналах и выступают на ведущих конференциях: \semibold{NeurIPS, ICML, ICLR, AISTATS, KDD, CLEF, LREC, WISE}.
  \subhead{Наши мероприятия}
  \begin{columns}[T]
    \begin{column}{0.7\textwidth}
      \begin{enumerate}
        \item Всероссийская научная конференция МФТИ --- секция проблем интеллектуального анализа данных, распознавания и прогнозирования
        \item Всероссийская конференция с международным участием «Математические методы распознавания образов»
        \item Международная конференция «Интеллектуализация обработки информации»
      \end{enumerate}
    \end{column}
    \begin{column}{0.26\textwidth}
      \centering
      \includegraphics[width=0.55\linewidth]{mipt.png}\\[0.2cm]
      \includegraphics[width=0.6\linewidth]{mmro.png}\\[0.2cm]
      \includegraphics[width=0.4\linewidth]{idp.png}
    \end{column}
  \end{columns}
\end{frame}

\begin{frame}{Выпуск 2026}
  \centering
  \vspace{0.2cm}
  \begin{tabular}{@{}c@{\hspace{0.4cm}}c@{}}
    \roundphoto{bachelors-2026.jpeg}{4.6cm} & \roundphoto{masters-2026.jpeg}{4.6cm}\\[0.1cm]
    \scriptsize\muted{Бакалавры} & \scriptsize\muted{Магистры}
  \end{tabular}
\end{frame}

\begin{frame}{Приём на кафедру}
  \begin{columns}[T]
    \begin{column}{0.62\textwidth}
      \begin{enumerate}
        \item Заполните анкету на странице \semibold{intsystems.github.io/admission}
        \item Выберите задачу из списка и решите её
        \item Расскажите о решении на собеседовании в формате презентации
        \item Расскажите, какая работа студента кафедры вам понравилась
        \item Дождитесь решения деканата
      \end{enumerate}
      \subhead{Даты}
      \footnotesize\admissiondates
    \end{column}
    \begin{column}{0.34\textwidth}
      \centering
      \qrlink[3.4cm]{https://intsystems.github.io/admission/}{intsystems.github.io/admission}
    \end{column}
  \end{columns}
\end{frame}

\coverframe[design-3]{%
  \covertitle{Ждём вас\\на кафедре!}{}%
  \node[anchor=south west] at ($(current page.south west)+(0.8cm,1.9cm)$) {%
    \qrlink{https://intsystems.github.io}{intsystems.github.io}\hspace{0.5cm}%
    \qrlink{https://t.me/IS_MIPT}{t.me/IS\_MIPT}\hspace{0.5cm}%
    \qrlink{https://www.youtube.com/c/MachineLearningPhystech}{YouTube}\hspace{0.5cm}%
    \qrlink{https://github.com/intsystems}{GitHub}%
  };
}

\end{document}
